% ============================================================================
% Fichier  : annexes/Ann_A_glossaire.tex
% Titre    : Glossaire
% Version  : 1.0 -- novembre 2026
% ============================================================================
% Définitions reprises des encadrés « Définition » des chapitres (termes
% marqués \termecle), abrégées ; le chapitre de référence est indiqué.
% ============================================================================

\chapter{Glossaire}
\label{ann:glossaire}

Ce glossaire reprend, sous une forme abrégée, les définitions données dans
les chapitres. Le numéro renvoie au chapitre où le terme est défini et
illustré~; c'est là qu'il faut lire la définition complète.

\begin{description}
    \item[Acceptation du risque] Décision argumentée, prise par l'instance
          qui en a l'autorité, de ne pas traiter davantage un risque
          résiduel, pour une durée et à des conditions données
          (Ch.~\ref{chap:iso27005}).
    \item[Appétence au risque] Niveau de risque qu'une organisation accepte
          pour atteindre ses objectifs~; fixée par l'organe de direction
          (Ch.~\ref{chap:grc}, \ref{chap:risque}).
    \item[Attributs de mesure] Étiquettes associées à chaque mesure de
          l'ISO/IEC~27002:2022 (type, propriétés, concepts, capacités,
          domaines), qui permettent de trier les mesures
          (Ch.~\ref{chap:iso27000}).
    \item[Audit] Processus méthodique, indépendant et documenté, qui obtient
          des preuves et les évalue au regard de critères
          (Ch.~\ref{chap:audit}).
    \item[Audit continu] Analyse répétée de toute une population de
          données ou de transactions pour y détecter les anomalies
          (Ch.~\ref{chap:emergents}).
    \item[Auditologie] Science appliquée de la réduction de l'incertitude
          sur les missions auditologiques (A0~; Ch.~\ref{chap:auditologie}).
    \item[Base commune de contrôles] Liste de contrôles propres à
          l'organisation, chacun relié aux référentiels qui l'exigent~;
          testée une fois, projetée sur plusieurs référentiels
          (Ch.~\ref{chap:referentiels}).
    \item[Base de preuve] Ce sur quoi repose une cote~: mesures internes,
          données externes, avis motivé, ou rien~; elle se cote elle-même
          \OUI{}, \NON{} ou \NEK{} (Ch.~\ref{chap:risque}).
    \item[Bien support] Composant sur lequel repose une valeur métier~:
          matériel, logiciel, réseau, personne, prestataire
          (Ch.~\ref{chap:risque}).
    \item[Certificat] Attestation, par un organisme tiers accrédité, qu'un
          système de management satisfait une norme, sur un périmètre et
          pour une durée donnés (Ch.~\ref{chap:prologue}).
    \item[Conclusion d'audit] Résultat de l'audit, tiré de l'ensemble des
          constats, au regard des objectifs de la mission
          (Ch.~\ref{chap:audit}).
    \item[Conformité] Satisfaction prouvée d'une exigence (norme, loi,
          contrat, politique) (Ch.~\ref{chap:prologue}, \ref{chap:grc}).
    \item[Constat d'audit] Résultat de l'évaluation des preuves au regard
          des critères~: conforme, non conforme, ou \NEK{}
          (Ch.~\ref{chap:audit}).
    \item[Contradictoire] Présentation des constats à l'audité, qui peut
          contester un fait, apporter une preuve ou expliquer une cause
          (Ch.~\ref{chap:audit}).
    \item[Contrôle continu de la conformité] Test automatique et répété
          d'un contrôle par une extraction comparée à une règle
          (Ch.~\ref{chap:grc-pratique}).
    \item[Couple SR/OV] Association d'une source de risque et d'un objectif
          visé, dans EBIOS Risk Manager (Ch.~\ref{chap:ebios}).
    \item[Criticité] Produit vraisemblance $\times$ impact, de 1 à 16~; elle
          classe, elle ne décide pas (Ch.~\ref{chap:risque}).
    \item[Critère d'audit] Exigence de référence~: norme, politique, loi,
          contrat, plan (Ch.~\ref{chap:audit}).
    \item[Critères d'acceptation] Règles, fixées à l'avance, qui disent quel
          résiduel peut être accepté, par qui, à quelles conditions
          (Ch.~\ref{chap:iso27005}).
    \item[Déclaration d'applicabilité] Document exigé par l'ISO/IEC~27001 qui
          recense les mesures nécessaires, justifie leur inclusion et
          l'exclusion des autres, et indique leur mise en œuvre
          (Ch.~\ref{chap:iso27000}).
    \item[Divergence de Kullback-Leibler] Mesure du changement de la
          distribution des réponses entre deux cycles
          (Ch.~\ref{chap:auditologie}).
    \item[Échantillon raisonné] Échantillon choisi selon les risques~; il
          trouve des défauts, il ne permet pas d'extrapoler un taux
          (Ch.~\ref{chap:audit}).
    \item[Entropie] Dispersion des réponses \OUI{}, \NON{}, \NEK{}, en bits,
          de 0 à $\log_2 3$~; elle se décompose en ignorance et désaccord
          (Ch.~\ref{chap:conclure}, \ref{chap:auditologie}).
    \item[Estimation calibrée] Minimum, valeur la plus probable et maximum,
          choisis pour encadrer la vraie valeur avec une confiance d'environ
          90~\% (Ch.~\ref{chap:fair}).
    \item[Événement redouté] Atteinte à une valeur métier, exprimée du point
          de vue de l'organisation (Ch.~\ref{chap:risque}).
    \item[Export figé] Copie d'une configuration ou d'une liste, extraite à
          une date et scellée par une empreinte (Ch.~\ref{chap:audit-technique}).
    \item[FAIR] Modèle de quantification du risque d'information~: fréquence
          des événements de perte et ampleur des pertes (Ch.~\ref{chap:fair}).
    \item[Fiche d'assertion] Description d'un contrôle par son objectif, son
          critère, sa question, son risque et sa conséquence
          (Ch.~\ref{chap:conclure}).
    \item[Gabarit canonique] Décomposition à huit niveaux~: domaine,
          sous-domaine, processus, activité, contrôle, test, preuve, trace
          (A0~; Ch.~\ref{chap:conclure}).
    \item[Gestion des risques] Activités coordonnées pour identifier,
          analyser, évaluer, traiter et suivre les risques
          (Ch.~\ref{chap:grc}, \ref{chap:iso27005}).
    \item[Gouvernance de la sécurité] Responsabilités et pratiques par
          lesquelles l'organe de direction fixe les objectifs, le niveau de
          risque acceptable, les rôles, et surveille les résultats
          (Ch.~\ref{chap:grc}).
    \item[GRC] Gouvernance, risque et conformité pilotés comme un seul
          dispositif, avec un référentiel commun (Ch.~\ref{chap:grc},
          \ref{chap:grc-pratique}).
    \item[Incertitude aléatoire] Part de l'incertitude qui tient au hasard
          des événements~; elle se décrit, elle ne se réduit pas
          (Ch.~\ref{chap:risque}).
    \item[Incertitude épistémique] Part de l'incertitude qui tient à ce que
          l'on ne sait pas~; elle se réduit par des preuves
          (Ch.~\ref{chap:risque}).
    \item[Indicateur de risque] Signal qui annonce un changement
          d'exposition, lu avec des seuils qui déclenchent une revue
          (Ch.~\ref{chap:grc-pratique}).
    \item[Matrice RACI] Tableau des rôles~: réalise, répond, consulté,
          informé~; un seul A par activité (Ch.~\ref{chap:grc}).
    \item[Matrice risques-contrôles] Tableau qui croise contrôles et
          risques, avec la dernière cote de chaque contrôle
          (Ch.~\ref{chap:grc-pratique}).
    \item[Menace] Cause potentielle d'un incident (Ch.~\ref{chap:risque}).
    \item[Mission auditologique] Quintuplet $\langle \mathcal{S},
          \mathcal{A}, \mathcal{I}, \mathcal{G}, \mathcal{T} \rangle$ (A0~;
          Ch.~\ref{chap:auditologie}).
    \item[Modèle des trois lignes] Répartition des rôles entre organe de
          gouvernance, première ligne, deuxième ligne et audit interne (IIA,
          2020~; Ch.~\ref{chap:grc}).
    \item[\NEK{}] Réponse «~connaissance insuffisante~»~: ni \OUI{}, ni
          \NON{}~; quatre conditions la rendent légitime (A1~;
          Ch.~\ref{chap:prologue}, \ref{chap:auditologie}).
    \item[Non-conformité majeure] Non-conformité qui affecte la capacité du
          système de management à atteindre ses résultats
          (Ch.~\ref{chap:smsi}).
    \item[Non-conformité mineure] Non-conformité qui n'affecte pas cette
          capacité (Ch.~\ref{chap:smsi}).
    \item[Objectif visé] Ce que la source de risque cherche à obtenir
          (Ch.~\ref{chap:risque}, \ref{chap:ebios}).
    \item[Organe de direction] Personne ou groupe qui répond en dernier
          ressort de l'organisation~: le conseil d'administration
          (Ch.~\ref{chap:grc}).
    \item[Périmètre de complaisance] Périmètre choisi pour faciliter la
          certification plutôt que pour couvrir les risques
          (Ch.~\ref{chap:prologue}).
    \item[Piste de progrès] Constat qui n'est pas une non-conformité, mais
          une occasion d'améliorer (Ch.~\ref{chap:smsi}).
    \item[Preuve d'audit] Information vérifiable et pertinente au regard des
          critères~; sa fiabilité se classe en cinq rangs
          (Ch.~\ref{chap:audit}).
    \item[Propriétaire de risque] Personne qui a la responsabilité et
          l'autorité de gérer un risque (Ch.~\ref{chap:grc}).
    \item[Rapport d'assurance] Rapport d'un auditeur indépendant sur les
          contrôles d'un prestataire (ISAE~3402, SOC~2)~; seul le type~2 dit
          s'ils ont fonctionné (Ch.~\ref{chap:referentiels}).
    \item[Registre des obligations] Liste des exigences qui s'imposent à
          l'organisation, avec leur source, leurs contrôles, leur preuve et
          leur statut (Ch.~\ref{chap:grc-pratique}).
    \item[Risque] Effet de l'incertitude sur les objectifs (ISO~31000~;
          Ch.~\ref{chap:risque}).
    \item[Risque cyber] Risque dont la source est une action sur les
          systèmes ou les données, et dont les conséquences touchent les
          activités (Ch.~\ref{chap:risque}).
    \item[Risque inhérent] Niveau du risque avant les mesures
          (Ch.~\ref{chap:risque}).
    \item[Risque résiduel] Niveau du risque compte tenu des seules mesures
          dont l'efficacité est établie (Ch.~\ref{chap:risque}).
    \item[Scénario opérationnel] Suite d'actions élémentaires (connaître,
          rentrer, trouver, exploiter) qui réalise un scénario stratégique
          (Ch.~\ref{chap:ebios}).
    \item[Scénario stratégique] Chemin d'attaque de haut niveau, d'un couple
          SR/OV à un événement redouté, éventuellement par l'écosystème
          (Ch.~\ref{chap:ebios}).
    \item[Sécurité effective] Capacité réelle de l'organisation à résister
          aux attaques et à s'en relever (Ch.~\ref{chap:prologue}).
    \item[Simulation de Monte-Carlo] Tirage d'un grand nombre d'années selon
          les estimations, pour obtenir la distribution des pertes
          (Ch.~\ref{chap:fair}).
    \item[SMSI] Système de management de la sécurité de l'information~:
          politiques, processus, rôles et contrôles sur un périmètre, en
          amélioration continue (Ch.~\ref{chap:prologue}, \ref{chap:smsi}).
    \item[Source de risque] Élément, personne ou groupe susceptible
          d'engendrer le risque (Ch.~\ref{chap:risque}).
    \item[Table de correspondance] Tableau qui relie les exigences de deux
          référentiels, en qualifiant chaque lien (Ch.~\ref{chap:referentiels}).
    \item[Taux de conformité ($\mu$)] Part des réponses \OUI{}
          (Ch.~\ref{chap:prologue}, \ref{chap:conclure}).
    \item[Taux d'incertitude ($U$)] Part des réponses \NEK{}
          (Ch.~\ref{chap:prologue}, \ref{chap:conclure}).
    \item[Valeur métier] Activité, processus ou information qui a de la
          valeur pour l'organisation (Ch.~\ref{chap:risque}).
    \item[Vulnérabilité] Faiblesse d'un bien support qu'une menace peut
          exploiter (Ch.~\ref{chap:risque}).
\end{description}
